\documentclass[11pt,a4paper]{article}

\usepackage[T1]{fontenc}
\usepackage[final]{microtype}
\usepackage[a4paper,margin=27mm,headheight=14pt]{geometry}
\usepackage{amsmath,amssymb,amsthm,mathtools}
\usepackage{newtxtext,newtxmath}
\usepackage{bm}
\usepackage{booktabs,longtable,array,tabularx}
\usepackage{enumitem}
\usepackage{xcolor}
\usepackage[most]{tcolorbox}
\usepackage{listings}
\usepackage{fancyhdr}
\usepackage{hyperref}
\usepackage[nameinlink,capitalise,noabbrev]{cleveref}

\definecolor{deepblue}{RGB}{25,63,112}
\definecolor{softblue}{RGB}{239,245,252}
\definecolor{softgray}{RGB}{247,247,247}
\definecolor{darkgray}{RGB}{72,72,72}

\hypersetup{
  colorlinks=true,
  linkcolor=deepblue,
  citecolor=deepblue,
  urlcolor=deepblue,
  pdftitle={Reversing x+W(x): Exact Reduction, All-Orders Asymptotics, and Logarithmic Transseries},
  pdfauthor={Research note prepared for the ProveIt project}
}

\pagestyle{fancy}
\fancyhf{}
\fancyhead[L]{\small Reversing \(x+W(x)\)}
\fancyhead[R]{\small Logarithmic transseries}
\fancyfoot[C]{\thepage}

\setlength{\parindent}{0pt}
\setlength{\parskip}{0.55em}
\setlist{nosep,leftmargin=2.0em}
\allowdisplaybreaks[3]
\numberwithin{equation}{section}

\newtheorem{theorem}{Theorem}[section]
\newtheorem{proposition}[theorem]{Proposition}
\newtheorem{lemma}[theorem]{Lemma}
\newtheorem{corollary}[theorem]{Corollary}
\theoremstyle{definition}
\newtheorem{definition}[theorem]{Definition}
\newtheorem{example}[theorem]{Example}
\theoremstyle{remark}
\newtheorem{remark}[theorem]{Remark}

\newtcolorbox{mainresultbox}{
  enhanced,
  breakable,
  colback=softblue,
  colframe=deepblue,
  boxrule=0.8pt,
  arc=1.2mm,
  left=3mm,right=3mm,top=2.5mm,bottom=2.5mm,
  title={Main result},
  fonttitle=\bfseries
}
\newtcolorbox{methodbox}[1][]{
  enhanced,
  breakable,
  colback=softgray,
  colframe=darkgray,
  boxrule=0.5pt,
  arc=1mm,
  left=3mm,right=3mm,top=2mm,bottom=2mm,
  #1
}

\lstdefinestyle{pythonstyle}{
  language=Python,
  basicstyle=\ttfamily\small,
  columns=fullflexible,
  keepspaces=true,
  frame=single,
  framerule=0.4pt,
  breaklines=true,
  showstringspaces=false,
  keywordstyle=\bfseries,
  commentstyle=\itshape,
  xleftmargin=0.5em,
  xrightmargin=0.5em
}

\newcommand{\e}{\mathrm e}
\newcommand{\dd}{\,\mathrm d}
\newcommand{\Ord}{\mathcal O}
\newcommand{\littleo}{\mathrm o}
\newcommand{\rW}{\mathcal W}
\newcommand{\Lone}{L_1}
\newcommand{\Ltwo}{L_2}
\newcommand{\Acoef}{A}
\newcommand{\Ccoef}{C}
\newcommand{\coeff}[2]{[\,#1^{#2}\,]}
\newcommand{\stirlingone}[2]{\genfrac{[}{]}{0pt}{}{#1}{#2}}
\newcommand{\invcomp}{\langle-1\rangle}
\newcommand{\Tfield}{\mathscr T}
\newcommand{\Cfield}{\mathscr C}
\newcommand{\val}{\nu}
\DeclareMathOperator{\Res}{Res}

\title{\textbf{Reversing \(x+W(x)\)}\\[0.35em]
\Large Exact \(r\)-Lambert Reduction, All-Orders Asymptotics,\\
and a General Calculus for Logarithmic Transseries}
\author{Research note prepared for the ProveIt project}
\date{2 September 2026}

\begin{document}
\maketitle

\begin{abstract}
Let \(W=W_0\) denote the real principal Lambert function and let
\(f(x)=x+W(x)\) on \([0,\infty)\).  We determine the inverse
\(g=f^{-1}\) at infinity in three mutually compatible forms.  First, an
exact change of variable shows that
\[
  g(z)=z-\rW_1(z),
\]
where \(\rW_r\) is the \(r\)-Lambert function, defined as the inverse of
\(y\mapsto y\e^y+ry\).  Second, if \(L=W(z)\) and
\(q=\e^{-L}=L/z\), then an analytic reversion problem at the ordinary
origin gives the convergent large-\(z\) expansion
\[
  g(z)=z-L+\sum_{n\ge1}\Acoef_n(L)q^n.
\]
We give a Lagrange coefficient formula, a triangular recurrence, the first
several rational coefficients, a joint expansion in \(q\) and \(1/L\), and
uniform remainder estimates.  Third, Lagrange--B\"urmann inversion at
infinity yields
\[
  g(z)=z+\sum_{n\ge1}\frac{(-1)^n}{n!}
       \frac{\dd^{n-1}}{\dd z^{n-1}}W(z)^n,
\]
which becomes a rational-coefficient expansion in powers of \(1/z\).
Substitution of the classical Stirling-number expansion of \(W\) produces
the full polynomial-logarithmic transseries in \(\log z\) and
\(\log\log z\).  A key point is that the correction to \(z-W(z)\) starts
at order \(W(z)/z\), beyond every fixed inverse power of \(\log z\).

The second half of the article extracts a reusable method for reversing
near-identity logarithmic transseries.  It treats slow coefficient fields,
valuation gain under differentiation, fixed-point and coefficient
recursions, Lagrange--B\"urmann inversion, Newton acceleration, domain
promotion, residual certification, and the family
\(x+\alpha W(x)\), as well as the more general perturbation
\(x+h(W(x))\).
\end{abstract}

\begin{mainresultbox}
For \(z\to+\infty\), put \(L=W(z)\).  Then
\begin{align*}
 f^{-1}(z)={}&z-L
 +\frac{L^2}{z(1+L)}
 +\frac{L^3(L^2-2)}{2z^2(1+L)^3}\\
 &+\frac{L^4(2L-1)(L^3-6L-6)}{6z^3(1+L)^5}\\
 &+\frac{L^5(6L^6-8L^5-69L^4-40L^3+96L^2+72L-24)}
 {24z^4(1+L)^7}
 +\Ord\!\left(\frac{L^5}{z^5}\right).
\end{align*}
More generally, after retaining the terms through \(z^{-N}\), the error is
\[
 f^{-1}(z)-G_N(z)
 \sim \frac{1}{N+1}\left(\frac{W(z)}{z}\right)^{N+1}.
\]
At the coarser inverse-logarithmic scale, with
\(\Lone=\log z\) and \(\Ltwo=\log\log z\),
\begin{align*}
 f^{-1}(z)\sim{}&z-\Lone+\Ltwo-\frac{\Ltwo}{\Lone}
 +\frac{\Ltwo(2-\Ltwo)}{2\Lone^2}\\
 &-\frac{\Ltwo(6-9\Ltwo+2\Ltwo^2)}{6\Lone^3}
 +\frac{\Ltwo(12-36\Ltwo+22\Ltwo^2-3\Ltwo^3)}{12\Lone^4}
 +\cdots.
\end{align*}
The first displayed expansion resolves the transseries sectors that the
second one cannot see.
\end{mainresultbox}

\tableofcontents
\newpage

\section{Introduction and orientation}

The Lambert function is defined by
\begin{equation}\label{eq:W-def}
  W(x)\e^{W(x)}=x.
\end{equation}
On the nonnegative real axis the principal branch is nonnegative, smooth,
and increasing.  The map
\begin{equation}\label{eq:f-def}
  f(x)=x+W(x),\qquad x\ge0,
\end{equation}
is therefore strictly increasing from \([0,\infty)\) onto itself, and its
inverse \(g=f^{-1}\) is globally defined there.

At first sight the inverse appears to be only a routine near-identity
reversion, because \(W(x)=\Ord(\log x)=\littleo(x)\).  The problem is more
instructive than that description suggests.  Three asymptotic scales occur:
\begin{enumerate}
\item the leading scale \(z\);
\item the slow logarithmic scale generated by \(\Lone=\log z\) and
      \(\Ltwo=\log\Lone\);
\item the much smaller scale \(q=W(z)/z=\e^{-W(z)}\), whose powers form
      separate transseries sectors.
\end{enumerate}
A calculation performed only in powers of \(1/\log z\) determines
\(z-W(z)\) to arbitrary inverse-logarithmic order, but it cannot detect the
first correction of size \(W(z)/z\).  This article keeps the scales
separate and then explains how to combine them safely.

The linked ProveIt draft collection develops algebra, composition, and
reversion for polynomial-logarithmic transseries~\cite{ProveItTransseries}.  The present note applies
that general viewpoint to one focused inverse problem and pushes the
special structure far enough to obtain an exact ordinary power-series
reversion, explicit rational coefficients, and a convergent sector
expansion.

\subsection{Notation and the three coordinate systems}

Throughout, \(W\) without a branch label means \(W_0\) on
\([0,\infty)\).  To avoid confusing a classical branch index with the
parameter of a generalized Lambert function, we use the following notation.

\begin{definition}[The \(r\)-Lambert function]
For a fixed real parameter \(r\), let \(\rW_r\) denote a local inverse of
\begin{equation}\label{eq:rW-def}
  y\longmapsto y\e^y+ry.
\end{equation}
When \(r=1\), this map is strictly increasing on the real line, so
\(\rW_1:\mathbb R\to\mathbb R\) is single-valued.  The notation
\(\rW_r\) is used here in place of the common notation \(W_r\), because
\(W_1\) also denotes a nonprincipal complex branch of the classical
Lambert function.  This generalized inverse is the \(r\)-Lambert function
studied by Mez\H{o} and Baricz~\cite{MezoBaricz2017}.
\end{definition}

We shall move among:
\[
 x \quad\xleftrightarrow{\ y=W(x)\ }\quad y,
 \qquad
 z \quad\xleftrightarrow{\ L=W(z)\ }\quad L,
 \qquad
 q=\e^{-L}=\frac{L}{z}.
\]
The \(y\)-coordinate gives the exact generalized-Lambert reduction.  The
\(L\)-coordinate rationalizes differentiation.  The \(q\)-coordinate
turns the exponentially small correction into an ordinary convergent power
series.

\section{Exact inversion and global real-variable properties}

\subsection[Reduction to the r-Lambert function]{Reduction to the \(r\)-Lambert function}

\begin{theorem}[Exact inverse]\label{thm:exact-inverse}
For every \(z\ge0\),
\begin{equation}\label{eq:exact-rW}
  g(z)=f^{-1}(z)=z-\rW_1(z)
       =\rW_1(z)\e^{\rW_1(z)}.
\end{equation}
Equivalently, if \(y=\rW_1(z)\), then \(g(z)=y\e^y\) and
\(W(g(z))=y\).
\end{theorem}

\begin{proof}
Put \(y=W(x)\).  Then \(x=y\e^y\), and the equation
\(z=x+W(x)\) becomes
\[
  z=y\e^y+y.
\]
By definition, \(y=\rW_1(z)\).  Since \(z=x+y\), we obtain
\(x=z-y=z-\rW_1(z)\).  The identity \(x=y\e^y\) gives the second form in
\eqref{eq:exact-rW}.  All quantities are nonnegative, so the principal real
branches are the relevant ones throughout.
\end{proof}

\paragraph{Lean status.}
Formal, in \path{LambertShiftInverse.lean}: \(\rW_1\) is \path{Fabius.shiftedLambertW}
(the inverse on \([0,\infty)\) of \path{Fabius.shiftedMulExp}, \(u\mapsto u\e^u+u\)), \(g\) is
\path{Fabius.lambertShiftInv}, and the three identities are
\path{Fabius.lambertShift_lambertShiftInv} (\(f(g(z))=z\)),
\path{Fabius.lambertShiftInv_eq_mul_exp} (\(g(z)=y\e^y\)) and
\path{Fabius.principalLambertW_lambertShiftInv} (\(W(g(z))=y\)); \(g\) is the two-sided
inverse of \(f\) on \([0,\infty)\) by \path{Fabius.lambertShiftInv_lambertShift} and
\path{Fabius.lambertShift_eq_iff}.  The bracket \eqref{eq:f-prime-bounds} is
\path{Fabius.one_le_deriv_lambertShift}, \path{Fabius.deriv_lambertShift_le_two}, and the
strict form of \eqref{eq:global-bracket} for \(z>0\) is \path{Fabius.lambertShiftInv_lt_self},
\path{Fabius.sub_principalLambertW_lt_lambertShiftInv}.  The derivative bracket
\eqref{eq:g-prime-bounds} is formal on the open half-line, in its strict form: \(g\) is
differentiable at every \(z>0\) by the inverse-function rule
(\path{Fabius.lambertShiftInv_hasDerivAt}, resting on the strict monotonicity and continuity of
\(g\), \path{Fabius.lambertShiftInv_strictMonoOn} and \path{Fabius.lambertShiftInv_continuousAt},
the latter read off the exact image \path{Fabius.lambertShiftInv_image_Ioi}), and
\path{Fabius.deriv_lambertShiftInv_bounds} is \(1/2<g'(z)<1\), both inequalities strict because
\(0<W'<1\) once \(W>0\) (\path{Fabius.inv_exp_mul_add_one_lt_one}).  The endpoint value
\(g'(0)=1/2\), a one-sided derivative, and the convexity of \(g\) are not formalized.

This exact formula identifies the special function involved, but it does
not by itself expose the full large-\(z\) hierarchy.  The remainder of the
article derives that hierarchy directly.

\subsection{Monotonicity, derivative bounds, and elementary brackets}

Differentiating \eqref{eq:W-def} gives
\begin{equation}\label{eq:W-prime}
 W'(x)=\frac{W(x)}{x(1+W(x))}
      =\frac{\e^{-W(x)}}{1+W(x)},
 \qquad x>0,
\end{equation}
with the continuous value \(W'(0)=1\).  Hence
\begin{equation}\label{eq:f-prime-bounds}
  1\le f'(x)=1+W'(x)\le2,
  \qquad x\ge0.
\end{equation}
This proves global invertibility and also yields
\begin{equation}\label{eq:g-prime-bounds}
  \frac12\le g'(z)<1,
  \qquad z\ge0.
\end{equation}
Since \(W''(x)<0\) for \(x>0\), the function \(f\) is concave and its
inverse \(g\) is convex.

The defining equation can also be written
\begin{equation}\label{eq:g-fixedpoint}
  g(z)=z-W(g(z)).
\end{equation}
Because \(0\le g(z)\le z\) and \(W\) is increasing, we obtain the useful
global bracket
\begin{equation}\label{eq:global-bracket}
  z-W(z)\le g(z)\le z.
\end{equation}
The lower approximation \(z-W(z)\) is already correct through every fixed
inverse-logarithmic order.  Its positive error is nevertheless of order
\(W(z)/z\), as will be proved exactly below.

\begin{proposition}[Residual-to-error certificate]\label{prop:residual-certificate}
Let \(\widetilde g(z)\ge0\) be any approximation and set
\begin{equation}\label{eq:residual-def}
  R(z)=\widetilde g(z)+W(\widetilde g(z))-z.
\end{equation}
Then
\begin{equation}\label{eq:residual-sandwich}
  \frac{|R(z)|}{2}\le |\widetilde g(z)-g(z)|\le |R(z)|.
\end{equation}
Moreover, \(R(z)\) and \(\widetilde g(z)-g(z)\) have the same sign.
\end{proposition}

\begin{proof}
Apply the mean-value theorem to \(f(\widetilde g)-f(g)=R\) and use
\eqref{eq:f-prime-bounds}.
\end{proof}

\paragraph{Lean status.}
The sandwich \eqref{eq:residual-sandwich} is formal for every \(z\ge0\) and every
\(\widetilde g(z)\ge0\): \path{Fabius.abs_sub_lambertShiftInv_le_abs_residual} and
\path{Fabius.abs_residual_le_two_mul_abs_sub_lambertShiftInv} in
\path{LambertShiftInverse.lean}, from the two-sided Lipschitz bracket of \(f\)
(\path{Fabius.abs_sub_le_abs_lambertShift_sub}, \path{Fabius.abs_lambertShift_sub_le_two_mul}),
itself the mean-value theorem with \eqref{eq:f-prime-bounds}.  The sign statement is
\path{Fabius.residual_pos_iff} and \path{Fabius.residual_neg_iff}.

This elementary certificate will later turn a formal composition check into
a rigorous numerical error bound.

\section{The natural small parameter and analytic reversion}

\subsection{Straightening the dominant Lambert behavior}

Let
\begin{equation}\label{eq:Lq-def}
  L=W(z),\qquad z=L\e^L,
  \qquad q=\e^{-L}=\frac{L}{z}.
\end{equation}
Write the generalized-Lambert solution in the form
\begin{equation}\label{eq:y-L-v}
  \rW_1(z)=L-v.
\end{equation}
By \cref{thm:exact-inverse},
\begin{equation}\label{eq:g-z-L-v}
  g(z)=z-L+v.
\end{equation}
Substitution of \eqref{eq:y-L-v} into
\((L-v)\e^{L-v}+L-v=L\e^L\), followed by division by \(\e^L\), gives
\[
  (L-v)(\e^{-v}+\e^{-L})=L.
\]
Thus the correction \(v\) is characterized exactly by
\begin{equation}\label{eq:B-L-equation}
  q=B_L(v):=\frac{L}{L-v}-\e^{-v}.
\end{equation}
The right side vanishes at \(v=0\) and
\[
  B_L'(0)=1+\frac1L>0.
\]
We have therefore reduced the large-variable problem to ordinary local
series inversion at the origin.

It is convenient to set
\begin{equation}\label{eq:u-def}
  u=\frac1L.
\end{equation}
Then
\begin{equation}\label{eq:Bvu}
  B(v,u)=\frac1{1-uv}-\e^{-v}
        =\sum_{k\ge1}b_k(u)v^k,
  \qquad
  b_k(u)=u^k-\frac{(-1)^k}{k!}.
\end{equation}
In particular, \(b_1(u)=1+u\).

\subsection{A convergent sector expansion}

\begin{theorem}[Joint analytic inverse]\label{thm:joint-analytic}
There are neighborhoods of \((q,u)=(0,0)\) and a unique analytic function
\(V(q,u)\) such that
\begin{equation}\label{eq:V-inverse}
  B(V(q,u),u)=q,
  \qquad V(0,u)=0.
\end{equation}
Consequently, for all sufficiently large positive \(z\),
\begin{equation}\label{eq:q-series-main}
  g(z)=z-L+V\!\left(\e^{-L},\frac1L\right)
      =z-L+\sum_{n\ge1}\Acoef_n(L)\e^{-nL},
  \qquad L=W(z),
\end{equation}
where the series converges.  Equivalently,
\begin{equation}\label{eq:q-series-L-over-z}
  g(z)=z-L+\sum_{n\ge1}\Acoef_n(L)
       \left(\frac{L}{z}\right)^n.
\end{equation}
The coefficients are analytic functions of \(u=1/L\) near \(u=0\).
\end{theorem}

\begin{proof}
The function
\[
  F(v,q,u)=B(v,u)-q
\]
is analytic near \((0,0,0)\), satisfies \(F(0,0,0)=0\), and has
\[
  \frac{\partial F}{\partial v}(0,0,0)=1.
\]
The analytic implicit-function theorem gives a unique analytic solution
\(v=V(q,u)\).  Along the positive-real path determined by
\(u=1/W(z)\) and \(q=\e^{-W(z)}\), both parameters tend to zero.  Hence
that path eventually enters the common domain of analyticity.  Substituting
into \eqref{eq:g-z-L-v} proves the result.
\end{proof}

This theorem is stronger than a formal asymptotic statement: the outer
sector series in powers of \(q=L/z\) is genuinely convergent for all
sufficiently large \(z\).  Each coefficient can then be expanded further
in powers of \(1/L\), producing the nested logarithmic transseries.

\subsection{Lagrange coefficient formula}

Since \(B_L(v)=b_1v+\Ord(v^2)\), ordinary Lagrange inversion gives the
following all-orders expression.

\begin{theorem}[Coefficient formula]\label{thm:A-Lagrange}
For \(n\ge1\),
\begin{align}
 \Acoef_n(L)
 &=\frac1n[v^{n-1}]
   \left(\frac{v}{B_L(v)}\right)^n
   \label{eq:A-coeff-extract}\\
 &=\frac1{n!}
   \left.\frac{\dd^{n-1}}{\dd v^{n-1}}
   \left(\frac{v}{\dfrac{L}{L-v}-\e^{-v}}\right)^n
   \right|_{v=0}.
   \label{eq:A-derivative}
\end{align}
Equivalently, in the variable \(u=1/L\), replace \(B_L(v)\) by
\(B(v,u)\) from \eqref{eq:Bvu}.
\end{theorem}

\begin{proof}
Apply the standard Lagrange inversion formula to \(q=B_L(v)\), or to
\(v=q\phi(v)\) with \(\phi(v)=v/B_L(v)\); see, for example,
Comtet~\cite{Comtet} or Flajolet and Sedgewick~\cite{FlajoletSedgewick}.
\end{proof}

\subsection{Triangular coefficient recursion}

The coefficient formula is compact, but a triangular recurrence is often
more efficient in exact symbolic arithmetic.  Write
\begin{equation}\label{eq:V-series}
  V(q,u)=\sum_{n\ge1}A_n(u)q^n,
  \qquad A_n(u)=\Acoef_n(1/u).
\end{equation}
Substituting into \(q=\sum_{k\ge1}b_k(u)V^k\) gives
\begin{equation}\label{eq:A1}
  A_1(u)=\frac1{1+u},
\end{equation}
and, for \(n\ge2\),
\begin{equation}\label{eq:A-recursion}
 A_n(u)=-\frac1{1+u}
 \sum_{k=2}^{n}b_k(u)
 \sum_{\substack{j_1+\cdots+j_k=n\\j_i\ge1}}
 A_{j_1}(u)\cdots A_{j_k}(u).
\end{equation}
The occurrence of \(A_n\) is linear and confined to the term
\(b_1A_n\); all other entries have lower indices.  This is the fundamental
triangularity behind near-identity reversion.

\begin{proposition}[Rational shape of the coefficients]\label{prop:A-shape}
For each \(n\ge1\), there is a polynomial \(R_n(L)\in\mathbb Z[L]\) of
degree \(2n-2\) such that
\begin{equation}\label{eq:A-shape}
  \Acoef_n(L)=\frac{L R_n(L)}{n!(1+L)^{2n-1}}.
\end{equation}
The leading coefficient of \(R_n\) is \((n-1)!\).  Consequently,
\begin{equation}\label{eq:A-limit}
  \Acoef_n(L)\longrightarrow\frac1n
  \qquad (L\to+\infty).
\end{equation}
\end{proposition}

\begin{proof}
Factor the linear coefficient from \(B(v,u)\):
\[
 B(v,u)=(1+u)v\left(1+
 \sum_{j\ge1}\frac{b_{j+1}(u)}{1+u}v^j\right).
\]
The coefficient form of Lagrange inversion therefore gives
\begin{align*}
 A_n(u)
 &=\frac{1}{n(1+u)^n}
 [v^{n-1}]
 \left(1+\sum_{j\ge1}
 \frac{b_{j+1}(u)}{1+u}v^j\right)^{-n}\\
 &=\frac1{n(1+u)^n}
 \sum_{m=0}^{n-1}(-1)^m
 \binom{n+m-1}{m}\frac{1}{(1+u)^m}
 [v^{n-1}]\left(\sum_{j\ge1}b_{j+1}(u)v^j\right)^m.
\end{align*}
After passage to the common denominator \((1+u)^{2n-1}\), every term
has numerator degree at most
\[
 (n-1-m)+\sum_{i=1}^{m}(j_i+1)
 =(n-1-m)+(n-1+m)=2n-2,
\]
where \(j_1+\cdots+j_m=n-1\).  It remains to check the asserted
integrality after multiplication by \(n!\).  Group the ordered
compositions by multiplicities \(c_j\), where
\(\sum_jc_j=m\) and \(\sum_jj c_j=n-1\).  Choosing from each
\(b_{j+1}(u)=u^{j+1}-(-1)^{j+1}/(j+1)!\) either its monomial part or its
constant part reduces the possible denominators to products of selected
factorials \((j+1)!\).  The potentially largest denominator is cleared by
\[
 (n-1)!\binom{n+m-1}{m}\frac{m!}{\prod_jc_j!}
 \prod_j\frac1{((j+1)!)^{c_j}}
 =\frac{(n+m-1)!}
 {\prod_jc_j!\,((j+1)!)^{c_j}},
\]
which is an integer: it counts set partitions of an
\((n+m-1)\)-element labelled set into \(c_j\) unlabeled blocks of size
\(j+1\).  This is the case in which every factor contributes its
constant part, and therefore it has the largest possible factorial
denominator.  If, for a given \(j\), only \(d_j\le c_j\) factors
contribute their constant parts, the grouped coefficient is the preceding
integer multiplied by
\(\binom{c_j}{d_j}((j+1)!)^{c_j-d_j}\), and similarly for all \(j\).
It is therefore integral as well.  Hence
\[
 A_n(u)=\frac{S_n(u)}{n!(1+u)^{2n-1}}
\]
for a polynomial \(S_n\in\mathbb Z[u]\) of degree at most \(2n-2\).
Substituting \(u=1/L\) yields
\[
 A_n(1/L)=\frac{L\,R_n(L)}{n!(1+L)^{2n-1}},
 \qquad R_n(L)=L^{2n-2}S_n(1/L)\in\mathbb Z[L].
\]
Finally, at \(u=0\) the inverse equation is
\[
 q=1-\e^{-v},
 \qquad v=-\log(1-q)=\sum_{n\ge1}\frac{q^n}{n},
\]
so \(A_n(0)=1/n\), equivalently \(S_n(0)=(n-1)!\).  Thus
\(R_n\) has degree exactly \(2n-2\) and leading coefficient
\((n-1)!\), completing the proof.
\end{proof}

\subsection{The first coefficients}

The first five coefficients are
\begin{align}
 \Acoef_1(L)&=\frac{L}{L+1},
 \label{eq:A1-L}\\
 \Acoef_2(L)&=\frac{L(L^2-2)}{2(L+1)^3},
 \label{eq:A2-L}\\
 \Acoef_3(L)&=\frac{L(2L-1)(L^3-6L-6)}{6(L+1)^5},
 \label{eq:A3-L}\\
 \Acoef_4(L)&=\frac{L}{24(L+1)^7}
 \bigl(6L^6-8L^5-69L^4-40L^3
       +96L^2+72L-24\bigr),
 \label{eq:A4-L}\\
 \Acoef_5(L)&=\frac{L}{120(L+1)^9}
 \bigl(24L^8-58L^7-428L^6-151L^5
 \notag\\
 &\hspace{34mm}{}+1320L^4+1380L^3-480L^2-720L+120\bigr).
 \label{eq:A5-L}
\end{align}
They can be generated either from \eqref{eq:A-derivative}, from
\eqref{eq:A-recursion}, or by the differential-operator method developed
in the next section.

\subsection[Resumming all q-sectors at fixed inverse powers of L]{Resumming all \(q\)-sectors at fixed inverse powers of \(L\)}

The joint analyticity in \((q,u)\) permits the opposite organization: first
expand in \(u=1/L\), while retaining the dependence on \(q\) exactly.  Put
\begin{equation}\label{eq:s-t-def}
  s=1-q,
  \qquad t=-\log(1-q).
\end{equation}
Solving \eqref{eq:Bvu} coefficient by coefficient in \(u\) gives
\begin{equation}\label{eq:joint-resummation}
 V(q,u)=t-\frac{ut}{s}
 +\frac{u^2t\bigl(2+(2q-1)t\bigr)}{2s^2}
 +\Ord(u^3),
\end{equation}
locally uniformly for \(q\) in a compact subset of \(|q|<1\).  Thus a
useful diagonally resummed approximation is
\begin{equation}\label{eq:resummed-approx}
 g(z)=z-L-\log(1-q)
 +\frac{\log(1-q)}{L(1-q)}
 +\Ord\!\left(\frac{-\log(1-q)}{L^2(1-q)^2}\right),
 \qquad q=\frac{L}{z}.
\end{equation}
At leading order in \(1/L\), all sector coefficients \(1/n\) sum to
\(-\log(1-q)\).  The ordinary expansion
\(\sum_n\Acoef_n(L)q^n\), however, is usually the cleaner form for
high-order exact computation.

\section{Lagrange--B\"urmann inversion at infinity}

The preceding section exploits a special exact coordinate.  A second
derivation starts directly from the near-identity map
\[
  f(x)=x+H(x),\qquad H(x)=W(x),
\]
and is the prototype for general logarithmic transseries reversal.

\subsection{The universal near-identity formula}

\begin{theorem}[Lagrange--B\"urmann formula at infinity]\label{thm:LB-infinity-special}
Formally, and asymptotically as \(z\to+\infty\),
\begin{equation}\label{eq:LB-infinity-special}
  g(z)=z+\sum_{n\ge1}\frac{(-1)^n}{n!}
  \frac{\dd^{n-1}}{\dd z^{n-1}}\bigl(W(z)^n\bigr).
\end{equation}
For every fixed power \(z^{-N}\), only finitely many summands contribute.
In the present problem the resulting expansion agrees term by term with the
convergent sector expansion \eqref{eq:q-series-L-over-z}.
\end{theorem}

\begin{proof}
Introduce a formal parameter \(\varepsilon\) and solve
\[
  K(\varepsilon)=-\varepsilon H(z+K(\varepsilon)),
  \qquad K(0)=0.
\]
Ordinary Lagrange inversion in \(\varepsilon\), with coefficients in the
differential algebra of functions of \(z\), gives
\[
  [\varepsilon^n]K(\varepsilon)
  =\frac{(-1)^n}{n!}D^{n-1}(H^n),
  \qquad D=\frac{\dd}{\dd z}.
\]
Setting \(\varepsilon=1\) yields \eqref{eq:LB-infinity-special}.  Here each
differentiation contributes one additional power of \(1/z\) in the
\(W\)-coordinate algebra, so the family is summable coefficientwise in
\(1/z\).  The sum satisfies the fixed-point equation
\(K=-H(z+K)\), hence is the inverse correction.  Agreement with
\eqref{eq:q-series-L-over-z} follows from uniqueness of the formal inverse,
or directly from \cref{prop:C-A-relation} below.
\end{proof}

Expanding the first few universal differential polynomials gives
\begin{align}
 (z+H)^{\langle-1\rangle}={}&z-H+HH'
 -H(H')^2-\frac12H^2H''
 \notag\\
 &+H(H')^3+\frac32H^2H'H''+\frac16H^3H'''
 +\cdots.
 \label{eq:universal-differential}
\end{align}
For a short perturbation \(H\), this formula often produces the first
blocks faster than direct composition.

\subsection[An exact differential operator in the W-coordinate]{An exact differential operator in the \(W\)-coordinate}

Because \(L=W(z)\),
\begin{equation}\label{eq:L-prime}
  \frac{\dd L}{\dd z}=\frac{L}{z(1+L)}.
\end{equation}
For a rational function \(R(L)\), define
\begin{equation}\label{eq:Tk-def}
  \mathcal D_k R
  :=\frac{L}{1+L}\frac{\dd R}{\dd L}-kR.
\end{equation}
Then
\begin{equation}\label{eq:D-operator-identity}
  \frac{\dd}{\dd z}\left(z^{-k}R(L)\right)
  =z^{-k-1}\mathcal D_kR(L).
\end{equation}
Thus differentiation raises the primary \(1/z\)-valuation by exactly one,
while all slow dependence remains in a rational coefficient.

\begin{theorem}[Coefficient operator formula]\label{thm:C-operator}
Define \(\Ccoef_n(L)\) for \(n\ge1\) by
\begin{equation}\label{eq:C-operator}
 \Ccoef_n(L)=\frac{(-1)^{n+1}}{(n+1)!}
 \bigl(\mathcal D_{n-1}\mathcal D_{n-2}\cdots
       \mathcal D_1\mathcal D_0\bigr)L^{n+1}.
\end{equation}
Then
\begin{equation}\label{eq:z-series-C}
  g(z)=z-L+\sum_{n\ge1}\frac{\Ccoef_n(L)}{z^n},
  \qquad L=W(z).
\end{equation}
Furthermore,
\begin{equation}\label{eq:C-shape}
  \Ccoef_n(L)=
  \frac{L^{n+1}R_n(L)}{n!(1+L)^{2n-1}},
\end{equation}
where the polynomial \(R_n\) is the same as in
\eqref{eq:A-shape}.
\end{theorem}

\begin{proof}
The term with index \(n+1\) in \eqref{eq:LB-infinity-special} is
\[
  \frac{(-1)^{n+1}}{(n+1)!}D^n(L^{n+1}).
\]
Repeated use of \eqref{eq:D-operator-identity}, beginning at \(k=0\),
gives \eqref{eq:C-operator} and shows that this term is exactly
\(z^{-n}\Ccoef_n(L)\).  On the other hand, the convergent sector
expansion \eqref{eq:q-series-L-over-z} writes the same inverse as
\[
 z-L+\sum_{n\ge1}z^{-n}L^n\Acoef_n(L).
\]
Uniqueness in the formal coefficient algebra \(\mathbb Q(L)[[z^{-1}]]\)
therefore gives
\(\Ccoef_n(L)=L^n\Acoef_n(L)\).  Substitution of
\eqref{eq:A-shape} proves \eqref{eq:C-shape} with the same polynomial
\(R_n\).
\end{proof}

\begin{proposition}[Equivalence of the two coefficient systems]
\label{prop:C-A-relation}
For every \(n\ge1\),
\begin{equation}\label{eq:C-A-relation}
  \Ccoef_n(L)=L^n\Acoef_n(L).
\end{equation}
Consequently,
\[
  \frac{\Ccoef_n(L)}{z^n}
  =\Acoef_n(L)\left(\frac{L}{z}\right)^n.
\]
\end{proposition}

\begin{proof}
Both \eqref{eq:q-series-L-over-z} and \eqref{eq:z-series-C} are expansions
of the unique inverse in powers of \(q=L/z\).  Comparing the coefficient
of \(q^n\) gives the identity.  It can also be verified immediately by
combining \eqref{eq:A-shape} and \eqref{eq:C-shape}.
\end{proof}

\subsection{Explicit rational-coefficient expansion}

The first five \(\Ccoef_n\) are
\begin{align}
 \Ccoef_1(L)&=\frac{L^2}{L+1},
 \label{eq:C1}\\
 \Ccoef_2(L)&=\frac{L^3(L^2-2)}{2(L+1)^3},
 \label{eq:C2}\\
 \Ccoef_3(L)&=\frac{L^4(2L-1)(L^3-6L-6)}{6(L+1)^5},
 \label{eq:C3}\\
 \Ccoef_4(L)&=\frac{L^5}{24(L+1)^7}
 \bigl(6L^6-8L^5-69L^4-40L^3
 +96L^2+72L-24\bigr),
 \label{eq:C4}\\
 \Ccoef_5(L)&=\frac{L^6}{120(L+1)^9}
 \bigl(24L^8-58L^7-428L^6-151L^5
 \notag\\[-0.2em]
 &\hspace{35mm}{}+1320L^4+1380L^3-480L^2-720L+120\bigr).
 \label{eq:C5}
\end{align}
Therefore
\begin{align}
 g(z)={}&z-L
 +\frac{L^2}{z(L+1)}
 +\frac{L^3(L^2-2)}{2z^2(L+1)^3}
 \notag\\
 &+\frac{L^4(2L-1)(L^3-6L-6)}{6z^3(L+1)^5}
 \notag\\
 &+\frac{L^5(6L^6-8L^5-69L^4-40L^3+96L^2+72L-24)}
 {24z^4(L+1)^7}
 \notag\\
 &+\frac{L^6(24L^8-58L^7-428L^6-151L^5+1320L^4+1380L^3
 -480L^2-720L+120)}
 {120z^5(L+1)^9}
 \notag\\
 &+\Ord\!\left(\frac{L^6}{z^6}\right),
 \qquad L=W(z).
 \label{eq:explicit-z-five}
\end{align}
The remainder notation in \eqref{eq:explicit-z-five} is justified by the
convergent \(q\)-series and the boundedness of \(\Acoef_n(L)\) for fixed
\(n\) as \(L\to\infty\).

\subsection[Large-L structure of the sector coefficients]{Large-\(L\) structure of the sector coefficients}

The leading behavior \(\Acoef_n(L)\to1/n\) implies
\begin{equation}\label{eq:C-leading}
  \Ccoef_n(L)\sim\frac{L^n}{n}.
\end{equation}
For example,
\begin{align}
 \Ccoef_1(L)&=L-1+\frac1L-\frac1{L^2}+\frac1{L^3}+\Ord(L^{-4}),
 \label{eq:C1-largeL}\\
 \Ccoef_2(L)&=\frac12L^2-\frac32L+2-\frac2L
 +\frac{3}{2L^2}+\Ord(L^{-3}),
 \label{eq:C2-largeL}\\
 \Ccoef_3(L)&=\frac13L^3-\frac{11}{6}L^2+\frac{23}{6}L
 -\frac{31}{6}+\frac{31}{6L}+\Ord(L^{-2}).
 \label{eq:C3-largeL}
\end{align}
Thus the \(n\)-th outer sector begins with
\((L/z)^n/n\), and summing only these leading diagonal terms recovers the
\(-\log(1-q)\) resummation in \eqref{eq:joint-resummation}.

\section[Expansion in log z and log log z]{Expansion in \(\log z\) and \(\log\log z\)}

\subsection{The classical Lambert expansion}

Set
\begin{equation}\label{eq:L1-L2}
  \Lone=\log z,
  \qquad \Ltwo=\log\Lone.
\end{equation}
Let \(\stirlingone{n}{k}\) denote an unsigned Stirling cycle number
of the first kind, and define
\begin{equation}\label{eq:Pn-def}
  P_n(s)=\sum_{m=1}^{n}(-1)^{n-m}
  \stirlingone{n}{n-m+1}\frac{s^m}{m!}.
\end{equation}
The standard large-argument expansion of the principal Lambert function is
\begin{equation}\label{eq:W-Stirling}
  W(z)\sim \Lone-\Ltwo+
  \sum_{n\ge1}\frac{P_n(\Ltwo)}{\Lone^n}.
\end{equation}
The first terms are
\begin{align}
 W(z)\sim{}&\Lone-\Ltwo+\frac{\Ltwo}{\Lone}
 +\frac{\Ltwo(\Ltwo-2)}{2\Lone^2}
 +\frac{\Ltwo(2\Ltwo^2-9\Ltwo+6)}{6\Lone^3}
 \notag\\
 &+\frac{\Ltwo(3\Ltwo^3-22\Ltwo^2+36\Ltwo-12)}{12\Lone^4}
 \notag\\
 &+\frac{\Ltwo(12\Ltwo^4-125\Ltwo^3+350\Ltwo^2-300\Ltwo+60)}
 {60\Lone^5}
 +\cdots.
 \label{eq:W-explicit}
\end{align}
This expansion is asymptotic in the Poincar\'e sense on the positive real
axis.  It is also obtainable by an ordinary convergent double-series
reversion in the variables \(1/\Lone\) and \(\Ltwo/\Lone\); see
Corless et al.~\cite{Corless1996} and DLMF \S4.13~\cite{DLMF}.

\subsection{The ordinary inverse-logarithmic sector}

Since \(g(z)=z-W(z)+\Ord(W(z)/z)\), the correction sectors in powers of
\(1/z\) are smaller than every fixed inverse power of \(\Lone\).  Thus,
for every fixed \(N\),
\begin{equation}\label{eq:g-log-sector-general}
 g(z)=z-\Lone+\Ltwo-
 \sum_{n=1}^{N}\frac{P_n(\Ltwo)}{\Lone^n}
 +\Ord\!\left(\frac{\Ltwo^{N+1}}{\Lone^{N+1}}\right).
\end{equation}
In explicit form,
\begin{align}
 g(z)\sim{}&z-\Lone+\Ltwo-\frac{\Ltwo}{\Lone}
 +\frac{\Ltwo(2-\Ltwo)}{2\Lone^2}
 \notag\\
 &-\frac{\Ltwo(6-9\Ltwo+2\Ltwo^2)}{6\Lone^3}
 +\frac{\Ltwo(12-36\Ltwo+22\Ltwo^2-3\Ltwo^3)}{12\Lone^4}
 \notag\\
 &-\frac{\Ltwo(12\Ltwo^4-125\Ltwo^3+350\Ltwo^2-300\Ltwo+60)}
 {60\Lone^5}
 +\cdots.
 \label{eq:g-log-explicit}
\end{align}
Equation \eqref{eq:g-log-explicit} is a valid answer at the ordinary
logarithmic scale.  It is not the complete transseries: any finite or
infinite formal manipulation confined to the \(z^0\) inverse-logarithmic
sector misses terms of order \(\Lone/z\), \(\Lone^2/z^2\), and so on.

\subsection{Why the first new sector is beyond all inverse-logarithmic orders}

For each fixed \(M\),
\begin{equation}\label{eq:beyond-all-log}
  \frac{\Lone/z}{\Lone^{-M}}=\frac{\Lone^{M+1}}{z}\longrightarrow0.
\end{equation}
Hence \(\Lone/z\) is beyond all algebraic orders in \(1/\Lone\).  The
first exact correction to \(z-W(z)\) is
\begin{equation}\label{eq:first-beyond-log}
  \frac{W(z)^2}{z(1+W(z))}
  \sim\frac{\log z}{z}.
\end{equation}
This explains a common apparent paradox: two approximations can agree to
all displayed inverse-logarithmic orders and still differ by a structured,
computable transseries sector.

A direct residual calculation makes the mechanism visible.  Let
\(G_0=z-L\).  Then
\[
  f(G_0)-z=W(z-L)-L.
\]
Taylor expansion around \(z\) gives
\[
  W(z-L)-L=-LW'(z)+\Ord(L^2W''(z))
  =-\frac{L^2}{z(1+L)}+\Ord\!\left(\frac{L^2}{z^2}\right).
\]
Since \(f'\sim1\), the correction needed to cancel this residual is
precisely the first term in \eqref{eq:explicit-z-five}.

\subsection{A rectangular polynomial-logarithmic transseries}

One may expand each rational function \(\Ccoef_n(W(z))\) in
\(\Lone^{-1}\) after fixing the outer sector \(z^{-n}\).  For example,
\begin{align}
 \Ccoef_1(W(z))={}&\Lone-\Ltwo-1
 +\frac{\Ltwo+1}{\Lone}
 +\frac{\Ltwo^2/2-1}{\Lone^2}
 \notag\\
 &+\frac{\Ltwo^3/3-\Ltwo^2/2-2\Ltwo+1}{\Lone^3}
 +\Ord\!\left(\frac{\Ltwo^4}{\Lone^4}\right),
 \label{eq:C1-log-expanded}\\
 \Ccoef_2(W(z))={}&\frac12\Lone^2-
 \left(\Ltwo+\frac32\right)\Lone
 +\frac12\Ltwo^2+\frac52\Ltwo+2
 \notag\\
 &-\frac{\Ltwo^2/2+5\Ltwo/2+2}{\Lone}
 +\Ord\!\left(\frac{\Ltwo^3}{\Lone^2}\right),
 \label{eq:C2-log-expanded}\\
 \Ccoef_3(W(z))={}&\frac13\Lone^3-
 \left(\Ltwo+\frac{11}{6}\right)\Lone^2
 +\left(\Ltwo^2+\frac{14}{3}\Ltwo+\frac{23}{6}\right)\Lone
 +\Ord(\Ltwo^3).
 \label{eq:C3-log-expanded}
\end{align}
Accordingly, in lexicographic transseries order,
\begin{align}
 g(z)\sim_{\mathrm{tr}}{}&
 z-\Lone+\Ltwo-\frac{\Ltwo}{\Lone}
 +\frac{\Ltwo(2-\Ltwo)}{2\Lone^2}-\cdots
 \notag\\
 &+\frac1z\left\{
 \Lone-\Ltwo-1+\frac{\Ltwo+1}{\Lone}
 +\frac{\Ltwo^2/2-1}{\Lone^2}+\cdots\right\}
 \notag\\
 &+\frac1{z^2}\left\{
 \frac12\Lone^2-\left(\Ltwo+\frac32\right)\Lone
 +\frac12\Ltwo^2+\frac52\Ltwo+2+\cdots\right\}
 \notag\\
 &+\frac1{z^3}\left\{
 \frac13\Lone^3-\left(\Ltwo+\frac{11}{6}\right)\Lone^2
 +\left(\Ltwo^2+\frac{14}{3}\Ltwo+\frac{23}{6}\right)\Lone
 +\cdots\right\}
 +\cdots.
 \label{eq:full-transseries-rectangular}
\end{align}
The symbol \(\sim_{\mathrm{tr}}\) emphasizes that the ordering is
hierarchical: first by the outer power \(z^{-n}\), then by the inverse-log
order within its coefficient.  A finite numerical approximation should
specify both cutoffs.

\section{Remainders, convergence, and certification}

\subsection{Uniform sector remainders}

Let
\begin{equation}\label{eq:GN-def}
  G_N(z)=z-L+\sum_{n=1}^{N}\Acoef_n(L)q^n,
  \qquad L=W(z),\quad q=\e^{-L}.
\end{equation}
The analytic inverse from \cref{thm:joint-analytic} immediately gives a
uniform remainder theorem.

\begin{theorem}[Sector remainder]\label{thm:sector-remainder}
For every fixed \(N\ge0\),
\begin{equation}\label{eq:sector-remainder-O}
  g(z)-G_N(z)=\Ord(q^{N+1})
  =\Ord\!\left(\frac{W(z)^{N+1}}{z^{N+1}}\right),
  \qquad z\to+\infty.
\end{equation}
More precisely,
\begin{equation}\label{eq:sector-remainder-asymp}
  g(z)-G_N(z)
  \sim \frac{q^{N+1}}{N+1}
  =\frac1{N+1}\left(\frac{W(z)}{z}\right)^{N+1}.
\end{equation}
In particular, the truncation error is positive for all sufficiently large
\(z\).
\end{theorem}

\begin{proof}
The convergent Taylor expansion of \(V(q,u)\) gives
\[
 V(q,u)-\sum_{n=1}^{N}A_n(u)q^n
 =A_{N+1}(u)q^{N+1}+\Ord(q^{N+2})
\]
uniformly for \(u\) near zero.  By \eqref{eq:A-limit},
\(A_{N+1}(u)\to1/(N+1)\).  Substitution of
\(u=1/L\) and \(q=L/z\) proves both claims.
\end{proof}

A quantitative Cauchy estimate is also available.  Choose a closed bidisc
\(|q|\le\rho\), \(|u|\le\eta\) lying inside the analytic domain of
\(V\), and let \(M=\max|V|\) on its distinguished boundary.  For
\(|q|<\rho\),
\begin{equation}\label{eq:Cauchy-remainder}
 \left|V(q,u)-\sum_{n=1}^{N}A_n(u)q^n\right|
 \le \frac{M(|q|/\rho)^{N+1}}{1-|q|/\rho}.
\end{equation}
This bound is conservative but completely explicit once a working bidisc
is chosen.

\subsection{Composition residuals}

A symbolic system need not know the analytic constants in
\eqref{eq:Cauchy-remainder}.  It can instead compute
\begin{equation}\label{eq:RN-def}
  R_N(z)=G_N(z)+W(G_N(z))-z.
\end{equation}
Formal cancellation through the retained sector gives
\[
  R_N(z)=\Ord(q^{N+1}).
\]
The residual certificate \cref{prop:residual-certificate} then supplies
an actual error enclosure:
\begin{equation}\label{eq:residual-error-GN}
  \frac{|R_N(z)|}{2}\le |G_N(z)-g(z)|\le |R_N(z)|.
\end{equation}
This is particularly valuable for computer algebra: construct the series
exactly, evaluate it numerically, evaluate one Lambert \(W\), and obtain a
rigorous a posteriori error bound without evaluating \(\rW_1\).

\subsection{Differentiated asymptotics}

The exact derivative identity
\begin{equation}\label{eq:g-derivative-exact}
  g'(z)=\frac1{1+W'(g(z))}
\end{equation}
provides an independent check.  Differentiating any truncation of
\eqref{eq:z-series-C} term by term and substituting it into
\eqref{eq:g-derivative-exact} must leave a residual beyond the target
order.  Because differentiation raises the \(1/z\)-valuation in the
\(W\)-coordinate field, the differentiated transseries is as well ordered
as the original one.

\section{Numerical verification}

For reference, the exact value in this section is computed by solving
\begin{equation}\label{eq:numeric-rW-equation}
  y\e^y+y=z
\end{equation}
for its unique nonnegative root and then setting \(g(z)=z-y\).  Let
\(G_N\) be the approximation \eqref{eq:GN-def}; thus \(N=0\) means
\(G_0=z-W(z)\).

\begin{table}[htbp]
\centering
\caption{Absolute error \(\lvert G_N(z)-g(z)\rvert\).}
\label{tab:numeric-errors}
\small
\begin{tabular}{rcccccc}
\toprule
\(z\) & \(N=0\) & \(N=1\) & \(N=2\) & \(N=3\) & \(N=4\) & \(N=5\)\\
\midrule
\(10\)
& \(1.12\times10^{-1}\) & \(1.05\times10^{-3}\) & \(2.99\times10^{-4}\)
& \(2.34\times10^{-5}\) & \(6.89\times10^{-7}\) & \(2.51\times10^{-7}\)\\
\(100\)
& \(2.64\times10^{-2}\) & \(2.19\times10^{-4}\) & \(9.50\times10^{-7}\)
& \(2.35\times10^{-8}\) & \(7.95\times10^{-10}\) & \(1.25\times10^{-11}\)\\
\(1000\)
& \(4.42\times10^{-3}\) & \(7.59\times10^{-6}\) & \(1.35\times10^{-8}\)
& \(1.72\times10^{-11}\) & \(1.10\times10^{-14}\) & \(1.81\times10^{-16}\)\\
\(10^6\)
& \(1.05\times10^{-5}\) & \(4.96\times10^{-11}\) & \(2.93\times10^{-16}\)
& \(1.84\times10^{-21}\) & \(1.16\times10^{-26}\) & \(7.16\times10^{-32}\)\\
\bottomrule
\end{tabular}
\end{table}

The corresponding exact values are
\begin{align*}
 g(10)&=8.366493829844153615806835\ldots,\\
 g(100)&=96.64072495463040645894733\ldots,\\
 g(1000)&=994.7548143481392806156898\ldots,\\
 g(10^6)&=999988.6166523780211244849\ldots.
\end{align*}
The data show the sector nature clearly.  Once \(z\) is moderately large,
each added block contributes roughly another factor \(q=W(z)/z\), in
accord with \eqref{eq:sector-remainder-asymp}.  At \(z=10\), where
\(q\) is not extremely small, the first several terms still improve the
approximation, although the early coefficients have not yet settled near
their limiting values \(1/n\).

\section{A general method for logarithmic-transseries reversal}

The special calculation above is an instance of a general and highly
algorithmic principle.  This section develops that principle in a form
suited to polynomial-logarithmic transseries at infinity.

\subsection{Slow coefficient fields and the primary valuation}

Let \(X\to+\infty\) and put \(t=X^{-1}\).  Introduce iterated logarithms
\begin{equation}\label{eq:iterated-logs}
  \ell_1=\log X,
  \qquad \ell_{j+1}=\log\ell_j.
\end{equation}
A depth-\(d\) slow coefficient field \(\Cfield_d\) may be taken to contain
rational functions and asymptotic Laurent series in these variables.  The
precise Hahn-field completion can be chosen to match the application; a
systematic framework is developed by van der Hoeven~\cite{vanDerHoeven}.
What matters for near-identity reversion here is closure under the slow
derivation
\begin{equation}\label{eq:slow-delta}
 \delta=
 \frac{\partial}{\partial\ell_1}
 +\frac1{\ell_1}\frac{\partial}{\partial\ell_2}
 +\frac1{\ell_1\ell_2}\frac{\partial}{\partial\ell_3}
 +\cdots
 +\frac1{\ell_1\cdots\ell_{d-1}}
  \frac{\partial}{\partial\ell_d}.
\end{equation}
Indeed, for \(a\in\Cfield_d\),
\begin{equation}\label{eq:D-slow}
  \frac{\dd a}{\dd X}=t\,\delta a.
\end{equation}
Define the outer block algebra
\begin{equation}\label{eq:Tfield-def}
  \Tfield_d=\Cfield_d[[t]].
\end{equation}
If \(a\in\Cfield_d\), then
\begin{equation}\label{eq:D-block-general}
  \frac{\dd}{\dd X}\bigl(t^na\bigr)
  =t^{n+1}(\delta a-na).
\end{equation}
Thus differentiation raises the primary \(t\)-valuation by one.  This
single property is the engine behind composition, fixed-point contraction,
and Lagrange summability.

The \(W\)-coordinate used earlier is an exact alternative slow field.  With
\(L=W(X)\), take \(\Cfield=\mathbb Q(L)\) and
\begin{equation}\label{eq:delta-W-field}
  \delta_W=\frac{L}{1+L}\frac{\dd}{\dd L}.
\end{equation}
Then \eqref{eq:D-block-general} becomes exactly
\eqref{eq:D-operator-identity}.  Choosing an adapted slow coordinate can
turn complicated logarithmic differentiation into rational algebra.

\subsection{Formal composition near the identity}

Let
\begin{equation}\label{eq:general-F-near-id}
  F(X)=X+H(X),
  \qquad H\in\Tfield_d.
\end{equation}
Seek an inverse
\begin{equation}\label{eq:general-G-near-id}
  G(X)=X+K(X),
  \qquad K\in\Tfield_d.
\end{equation}
Taylor composition is defined by
\begin{equation}\label{eq:Taylor-composition}
  H(X+K)=\sum_{j\ge0}\frac{H^{(j)}(X)}{j!}K^j.
\end{equation}
Because \(H^{(j)}\in t^j\Tfield_d\), only finitely many Taylor terms can
contribute to any prescribed outer order.  Thus the composition is
well-defined in the \(t\)-adic topology even when \(K\) begins at outer
order zero.

\begin{theorem}[Near-identity inverse theorem]\label{thm:general-near-id}
For every \(F=X+H\) with \(H\in\Tfield_d\), there is a unique
\(G=X+K\), \(K\in\Tfield_d\), such that
\begin{equation}\label{eq:two-sided-general}
  F\circ G=G\circ F=X.
\end{equation}
The correction is the unique fixed point of
\begin{equation}\label{eq:general-fixedpoint}
  K=-H(X+K).
\end{equation}
\end{theorem}

\begin{proof}
Let \(\Phi(K)=-H(X+K)\).  For two corrections \(K_1,K_2\), Taylor's
difference formula gives
\[
 H(X+K_1)-H(X+K_2)
 =\sum_{j\ge1}\frac{H^{(j)}(X+K_2)}{j!}(K_1-K_2)^j.
\]
The first derivative lies in \(t\Tfield_d\), and the \(j\)-th derivative
lies in \(t^j\Tfield_d\).  Therefore
\begin{equation}\label{eq:valuation-gain}
 \val_t\bigl(\Phi(K_1)-\Phi(K_2)\bigr)
 \ge \val_t(K_1-K_2)+1.
\end{equation}
Starting with \(K_0=0\) and iterating \(K_{m+1}=\Phi(K_m)\) produces a
Cauchy sequence in the complete \(t\)-adic algebra.  Its limit is the
unique right inverse correction.  A right inverse tangent to the identity
is also a left inverse: if \(J=G\circ F=X+R\), then associativity gives
\(J\circ J=J\).  If \(R\ne0\), the leading nonzero block of
\(J\circ J-X\) is twice that of \(R\), a contradiction in characteristic
zero.  Hence \(R=0\).
\end{proof}

\subsection{General Lagrange--B\"urmann formula}

\begin{theorem}[Formal reversal formula]\label{thm:general-LB}
Under the hypotheses of \cref{thm:general-near-id},
\begin{equation}\label{eq:general-LB}
  F^{\langle-1\rangle}(X)
  =X+\sum_{n\ge1}\frac{(-1)^n}{n!}
   D^{n-1}\bigl(H(X)^n\bigr),
  \qquad D=\frac{\dd}{\dd X}.
\end{equation}
The family is summable in the primary \(t\)-valuation.
\end{theorem}

\begin{proof}
Repeat the auxiliary-parameter proof of
\cref{thm:LB-infinity-special}.  Since \(H^n\) has nonnegative primary
valuation and \(D^{n-1}\) raises that valuation by at least \(n-1\), the
\(n\)-th summand begins at order \(t^{n-1}\).  Therefore only finitely
many summands affect a fixed block, and setting the auxiliary parameter to
one is formally legitimate.  The sum satisfies the fixed-point equation
and is identified by uniqueness.
\end{proof}

This is the inverse-at-infinity form of the classical Lagrange--B\"urmann
formula; standard analytic treatments may be found in
Henrici~\cite{Henrici}, while the formal coefficientwise version used here
is developed in the linked ProveIt drafts~\cite{ProveItTransseries}.

\begin{corollary}[Universal first blocks]\label{cor:universal-blocks}
The inverse begins as
\begin{align}
 F^{\langle-1\rangle}={}&X-H+HH'
 -H(H')^2-\frac12H^2H''
 \notag\\
 &+H(H')^3+\frac32H^2H'H''+\frac16H^3H'''
 +\cdots.
 \label{eq:universal-blocks-general}
\end{align}
\end{corollary}

\subsection{Direct coefficient recursion}

For symbolic work with many blocks, it is often better not to expand the
powers \(H^n\).  Suppose an approximation \(G_N\) has been determined
through outer order \(t^N\).  Introduce one unknown block
\(a(\ell_1,\ldots,\ell_d)t^{N+1}\), compose, and set the first residual
block to zero.  Near the identity the new block appears linearly with
coefficient one, so the equation is triangular.

\begin{methodbox}[title={Coefficient-by-coefficient reversal}]
\begin{enumerate}
\item Choose the transmonomial hierarchy and a coefficient field closed
      under the needed algebra and derivation.
\item Normalize the dominant part of \(F\).  If it is not \(X\), invert
      the leading monomial first.
\item Insert \(G_N=X+K_N\) into \(F\circ G_N-X\), retaining only the
      first unresolved block.
\item Add one unknown block to \(K_N\), solve the resulting linear
      coefficient equation, and truncate immediately.
\item Continue to the target order, then check both
      \(F\circ G_N-X\) and \(G_N\circ F-X\).
\end{enumerate}
\end{methodbox}

In the special \(q\)-reversion of \cref{eq:B-L-equation}, this procedure is
exactly the recurrence \eqref{eq:A-recursion}.

\subsection{Newton acceleration}

If \(G\) is an approximate inverse and
\begin{equation}\label{eq:general-residual}
  R=F\circ G-X,
\end{equation}
linearization gives the formal Newton update
\begin{equation}\label{eq:general-newton}
  G_{\mathrm{new}}
  =G-\frac{F\circ G-X}{F'\circ G}.
\end{equation}
Because \(F'=1+t\Tfield_d\), the denominator is a unit.  After the scale
and branch have been normalized correctly, Newton iteration approximately
doubles the number of valid outer blocks at each step.  A robust
implementation computes two or three blocks by the triangular method,
then switches to Newton.

\subsection{From a formal inverse to an actual asymptotic inverse}

Formal correctness becomes analytic correctness under mild stability
conditions; this residual-based passage is compatible with the standard
Poincar\'e asymptotic viewpoint~\cite{Olver}.  Suppose an actual function \(F\) admits the differentiated
transseries used in the calculation and that \(F'(x)\ge c>0\) eventually.
If a truncation \(G_N\) has composition residual
\begin{equation}\label{eq:formal-to-actual-residual}
  F(G_N(X))-X=\Ord(M_N(X)),
\end{equation}
then the mean-value theorem gives
\begin{equation}\label{eq:formal-to-actual-error}
  G_N(X)-F^{-1}(X)=\Ord(M_N(X)).
\end{equation}
Thus residual control is the natural bridge between formal transseries and
Poincare asymptotics.  In the present problem \(c=1\), and
\cref{prop:residual-certificate} gives a two-sided numerical enclosure.

\section{Scale selection, domain promotion, and common failure modes}

The formal theorem assumes that the leading map is the identity and that
all coefficients live in a differential field stable under the required
operations.  Most practical failures of transseries reversal occur before
the coefficient recursion begins: the dominant scale is normalized
incorrectly, an iterated logarithm is omitted, or a composition leaves the
chosen coefficient domain.

\subsection{Invert the dominant monomial first}

Suppose
\begin{equation}\label{eq:dominant-power-log}
  Y\sim cX^a(\log X)^b,
  \qquad ca\ne0.
\end{equation}
If \(b=0\), the leading inverse is \((Y/c)^{1/a}\).  If \(b\ne0\), put
\(s=\log X\).  Then
\[
  \left(\frac{Y}{c}\right)^{1/b}=s\e^{(a/b)s},
\]
and therefore
\begin{equation}\label{eq:power-log-leading-inverse}
  X=\exp\!\left[
  \frac ba W_k\!\left(
  \frac ab\left(\frac{Y}{c}\right)^{1/b}
  \right)\right].
\end{equation}
The branch \(k\) is selected by the real interval or complex sector under
study.  Expansion of \(W_k\) introduces \(\log Y\) and
\(\log\log Y\); thus a second logarithmic level is forced whenever the
leading monomial contains a nonzero power of \(\log X\).

After the leading map \(F_0\) is inverted, factor
\begin{equation}\label{eq:dominant-factorization}
  F=F_0\circ\Phi,
  \qquad \Phi=X+\text{smaller terms}.
\end{equation}
Then
\begin{equation}\label{eq:dominant-factorization-inverse}
  F^{\langle-1\rangle}
  =\Phi^{\langle-1\rangle}\circ F_0^{\langle-1\rangle},
\end{equation}
and the near-identity machinery applies only to \(\Phi\).

\subsection{Generator-first composition}

Before expanding an outer transseries at an inner argument \(G\), compute
the images of its generators.  For example, if
\begin{equation}\label{eq:general-inner-leading}
  G(X)=cX^\beta(\log X)^\gamma(1+U(X)),
  \qquad U\prec1,
\end{equation}
then
\begin{align}
  G^{-1}&=c^{-1}X^{-\beta}(\log X)^{-\gamma}(1+U)^{-1},
  \label{eq:generator-inverse-image}\\
  \log G&=\beta\log X+\gamma\log\log X+\log c+\log(1+U).
  \label{eq:generator-log-image}
\end{align}
The term \(\gamma\log\log X\) is not optional.  If the coefficient domain
contains only polynomials in \(\log X\), composition with
\((\log X)^\gamma\) generally forces promotion to a depth-two logarithmic
field.

\subsection{A practical promotion chain}

A safe symbolic implementation can use the explicit chain
\begin{equation}\label{eq:promotion-chain}
 \mathbb Q[\ell_1]
 \subset\mathbb Q(\ell_1)
 \subset\mathbb Q((\ell_1^{-1}))
 \subset\mathbb Q((\ell_1^{-1}))[\ell_2]
 \subset\cdots.
\end{equation}
Promotion is triggered by a concrete operation:
\begin{itemize}
\item division by a nonconstant polynomial coefficient requires a rational
      function field;
\item expansion for large \(\ell_1\) requires Laurent series in
      \(\ell_1^{-1}\);
\item inversion of a leading factor containing \((\log X)^b\) introduces
      \(\log\log X\);
\item composition with more complicated leading monomials may introduce
      additional iterated logarithms or real powers.
\end{itemize}
Every promotion changes the support order and the meaning of truncation.  It
should be recorded as part of the result, not hidden as a simplification.

\subsection{Failure modes}

The following errors are especially common.

\begin{description}[style=nextline,leftmargin=2em,labelindent=0pt]
\item[Using only the inverse-logarithmic sector.]
This gives \(z-W(z)\) correctly to arbitrary inverse-log order but misses
the \(W(z)/z\) sectors.  The remedy is to introduce the outer parameter
\(q=W(z)/z\) or work in \(\Cfield[[z^{-1}]]\).

\item[Expanding before choosing the correct slow coordinate.]
An expression that is rational in \(L=W(z)\) can become very large after
premature substitution of the \((\Lone,\Ltwo)\) series.  Compute exact
outer coefficients in \(L\) first and expand them in inverse logarithms
only at the end.

\item[Ignoring branch information.]
The positive real problem has no branch ambiguity, but complex extensions
do.  Classical branch indices, generalized-Lambert parameters, and chosen
logarithm branches must be kept distinct.

\item[Truncating only at the end.]
Taylor composition, powers, and rational simplification can create severe
expression swell.  Truncate after every algebraic operation and cache
repeated powers.

\item[Checking only the displayed coefficients.]
A plausible pattern is not a proof.  Verify a left composition residual, a
right composition residual, and, when possible, the differentiated inverse
identity.
\end{description}

\section{Parametric and functional extensions}

\subsection[The family x+alpha W(x)]{The family \(x+\alpha W(x)\)}

Let
\begin{equation}\label{eq:f-alpha}
  f_\alpha(x)=x+\alpha W(x),
  \qquad \alpha\ge0.
\end{equation}
With \(y=W(x)\), the inversion equation is
\begin{equation}\label{eq:f-alpha-y}
  y\e^y+\alpha y=z.
\end{equation}
Thus the exact inverse is
\begin{equation}\label{eq:f-alpha-exact}
  f_\alpha^{-1}(z)=z-\alpha\rW_\alpha(z).
\end{equation}
Set again \(L=W(z)\), \(q=\e^{-L}\), and write
\(\rW_\alpha(z)=L-v\).  The same calculation as before gives
\begin{equation}\label{eq:alpha-B-equation}
  B_L(v)=\alpha q.
\end{equation}
Therefore the coefficient functions \(\Acoef_n(L)\) are universal:
\begin{equation}\label{eq:f-alpha-q-expansion}
 f_\alpha^{-1}(z)
 =z-\alpha L
 +\alpha\sum_{n\ge1}\Acoef_n(L)(\alpha q)^n.
\end{equation}
Equivalently,
\begin{equation}\label{eq:f-alpha-z-expansion}
 f_\alpha^{-1}(z)
 =z-\alpha L
 +\sum_{n\ge1}\alpha^{n+1}
   \frac{\Ccoef_n(L)}{z^n}.
\end{equation}
For every fixed \(N\), the remainder after \(n=N\) is
\begin{equation}\label{eq:f-alpha-error}
 \sim \frac{\alpha^{N+2}}{N+1}
 \left(\frac{L}{z}\right)^{N+1}.
\end{equation}
The original problem is the case \(\alpha=1\).

\subsection[Perturbations of the form x+h(W(x))]{Perturbations of the form \(x+h(W(x))\)}

Consider
\begin{equation}\label{eq:f-h}
  F_h(x)=x+h(W(x)),
\end{equation}
where \(h(y)\) has subexponential growth as \(y\to+\infty\) and is regular
in the relevant slow coefficient field.  Put \(y=W(x)\), so
\begin{equation}\label{eq:f-h-y}
  z=y\e^y+h(y).
\end{equation}
Let \(L=W(z)\), \(q=\e^{-L}\), and \(y=L-v\).  Dividing
\eqref{eq:f-h-y} by \(\e^L\) and rearranging gives the exact implicit
relation
\begin{equation}\label{eq:f-h-B}
  B_L(v)=q\,\frac{h(L-v)}{L-v},
  \qquad
  B_L(v)=\frac{L}{L-v}-\e^{-v}.
\end{equation}
Once \(v\) is found as a \(q\)-series, the inverse is
\begin{equation}\label{eq:f-h-inverse}
  F_h^{-1}(z)=z-h(L-v).
\end{equation}
The leading correction follows immediately.  Since
\(B_L'(0)=(L+1)/L\),
\begin{equation}\label{eq:f-h-v-leading}
  v=\frac{h(L)}{L+1}q+\Ord(q^2),
\end{equation}
and hence
\begin{equation}\label{eq:f-h-leading-inverse}
  F_h^{-1}(z)
  =z-h(L)+\frac{h(L)h'(L)}{L+1}q+\Ord(q^2).
\end{equation}
For \(h(L)=\alpha L\), this reduces to
\eqref{eq:f-alpha-q-expansion}.  For a polynomial or polynomial-logarithmic
\(h\), every higher coefficient is obtained by the same triangular
composition procedure and remains in the chosen slow field.

\subsection{Ordinary near-identity perturbations}

For a general
\begin{equation}\label{eq:F-general-H}
  F(X)=X+H(X),
  \qquad H(X)=\littleo(X),
\end{equation}
the direct formula \eqref{eq:general-LB} is usually sufficient.  The
special transform through \(W\) is advantageous only when the perturbation
is naturally expressed in \(W(X)\), because then the exact differential
relation \eqref{eq:delta-W-field} and the identity
\(\e^{-W(X)}=W(X)/X\) expose the correct transseries sectors without
expanding nested logarithms prematurely.

\section{Algorithms and symbolic implementation}

\subsection[Operator generation of the 1/z coefficients]{Operator generation of the \(1/z\) coefficients}

The following short SymPy-style routine implements
\eqref{eq:C-operator}.  It keeps every coefficient as an exact rational
function of \(L\).

\begin{lstlisting}[style=pythonstyle,caption={Exact generation of the coefficient \(C_n(L)\).},label={lst:C-generator}]
from sympy import symbols, diff, factorial, factor

L = symbols("L")

def D_block(k, R):
    """Return D_k R = L/(1+L) R' - k R."""
    return L/(1 + L) * diff(R, L) - k*R

def C(n):
    """Coefficient of z**(-n) in the inverse of x + W(x)."""
    R = L**(n + 1)
    for k in range(n):              # D_0, D_1, ..., D_{n-1}
        R = D_block(k, R)
    return factor((-1)**(n + 1) * R / factorial(n + 1))

for n in range(1, 7):
    print(n, C(n))
\end{lstlisting}

The approximation itself is
\begin{lstlisting}[style=pythonstyle,caption={Evaluation of a sector truncation.}]
from mpmath import mp

mp.dps = 80

def inverse_approx(z, N, C_functions):
    Lval = mp.lambertw(z)           # principal real branch
    ans = z - Lval
    for n in range(1, N + 1):
        ans += C_functions[n](Lval) / z**n
    return ans
\end{lstlisting}

\subsection[Direct q-series reversion]{Direct \(q\)-series reversion}

A generic coefficient engine for \eqref{eq:Bvu} can instead operate on
truncated ordinary power series.

\begin{lstlisting}[style=pythonstyle,caption={Triangular reversion of \(q=B(v,u)\).}]
from sympy import symbols, exp, series, solve, expand

q, u, v = symbols("q u v")
B = 1/(1 - u*v) - exp(-v)
V = 0

for n in range(1, N + 1):
    an = symbols(f"a{n}")
    trial = V + an*q**n
    residual = series(B.subs(v, trial) - q, q, 0, n + 1).removeO()
    equation = expand(residual).coeff(q, n)
    an_value = solve(equation, an)[0]
    V = expand(V + an_value*q**n)
\end{lstlisting}

This approach is preferable when the implicit relation is modified, as in
\eqref{eq:f-h-B}.  The operator method is preferable for the original
near-identity problem because it avoids repeated series composition.

\subsection{Hybrid workflow for a general logarithmic transseries}

\begin{methodbox}[title={Recommended reversal workflow}]
\begin{enumerate}
\item Determine the dominant monomial and invert it exactly or
      asymptotically.  Fix real or complex branches at this stage.
\item Compute the images of \(X^{-1}\), \(\log X\), and all deeper
      generators under the proposed leading inverse.  Promote the
      coefficient field when a new scale appears.
\item Normalize to \(F=X+H\) in a complete outer block algebra.
\item Compute two or three blocks by direct coefficient cancellation.
      This tests the scale choice and exposes any missing generators.
\item Switch to Newton reversion when many blocks are required, doubling
      the working cutoff after each update.
\item Expand adapted slow coordinates, such as \(L=W(X)\), into ordinary
      iterated logarithms only after the exact outer coefficients are known.
\item Verify left, right, and derivative residuals.  For actual functions,
      convert the composition residual into an error bound using a lower
      bound for \(|F'|\).
\end{enumerate}
\end{methodbox}

\subsection{Complexity and expression control}

Exact arithmetic is strongly recommended.  Coefficients such as
\eqref{eq:C4} and \eqref{eq:C5} contain large cancellations that floating
point arithmetic can obscure.  Expression growth is controlled by:
\begin{itemize}
\item retaining factored rational functions until coefficients must be
      compared;
\item using Horner evaluation for slow-variable polynomials;
\item caching powers and derivatives;
\item truncating after every multiplication, composition, and division;
\item separating algebraic simplification from branch-dependent analytic
      simplification;
\item representing the two cutoffs explicitly: outer sector order and
      inner inverse-log order.
\end{itemize}

\section{Conclusions}

The inverse of \(x+W(x)\) has more structure than a single logarithmic
asymptotic series suggests.  The exact substitution \(y=W(x)\) identifies
it as
\[
  f^{-1}(z)=z-\rW_1(z).
\]
The further substitution \(L=W(z)\), followed by
\(\rW_1(z)=L-v\), converts the large-\(z\) problem into the local analytic
reversion
\[
  \frac{L}{L-v}-\e^{-v}=\e^{-L}.
\]
This yields a convergent expansion in the natural sector parameter
\(q=\e^{-L}=L/z\), with rational coefficient functions and a simple
Lagrange formula.  The same coefficients arise from Lagrange--B\"urmann
inversion at infinity through the differential operators
\(
\mathcal D_k=L(1+L)^{-1}\dd/\dd L-k
\).

The ordinary inverse-logarithmic expansion is obtained by replacing
\(L=W(z)\) by its Stirling-number series in \(\log z\) and
\(\log\log z\).  That sector alone is incomplete: the first correction to
\(z-W(z)\) is of order \(W(z)/z\), smaller than every fixed inverse power
of \(\log z\).  Keeping the outer \(1/z\) sectors separate from the slow
coefficient expansions resolves the hierarchy cleanly.

The reusable lesson is methodological.  Series reversal of logarithmic
transseries should proceed by dominant-scale normalization, generator-first
composition, work in a slow differential coefficient field, and
near-identity inversion by triangular recursion, Lagrange--B\"urmann, or
Newton iteration.  Composition residuals are not merely formal checks;
under a derivative lower bound they provide rigorous error certificates.

\appendix

\section{Additional exact coefficients}

With
\begin{equation}\label{eq:appendix-An-shape}
  \Acoef_n(L)=\frac{L R_n(L)}{n!(1+L)^{2n-1}},
  \qquad
  \Ccoef_n(L)=L^n\Acoef_n(L),
\end{equation}
the first six numerator polynomials are
\begin{align*}
 R_1(L)={}&1,\\
 R_2(L)={}&L^2-2,\\
 R_3(L)={}&2L^4-L^3-12L^2-6L+6,\\
 R_4(L)={}&6L^6-8L^5-69L^4-40L^3+96L^2+72L-24,\\
 R_5(L)={}&24L^8-58L^7-428L^6-151L^5+1320L^4+1380L^3
            -480L^2-720L+120,\\
 R_6(L)={}&120L^{10}-444L^9-2920L^8+424L^7+16577L^6
            +17892L^5\\
          &\quad{}-13530L^4-26400L^3-1800L^2+7200L-720.
\end{align*}
Thus
\begin{equation}\label{eq:A6}
 \Acoef_6(L)=
 \frac{L R_6(L)}{720(1+L)^{11}},
 \qquad
 \Ccoef_6(L)=
 \frac{L^7 R_6(L)}{720(1+L)^{11}}.
\end{equation}
The highest-degree coefficient of \(R_n\) is \((n-1)!\), while the
constant term for the displayed cases is \((-1)^{n-1}n!\).  The operator
formula \eqref{eq:C-operator} is preferable to storing these expanded
polynomials once \(n\) becomes large.

For comparison, the same coefficients in \(u=1/L\) begin as
\begin{align*}
 A_1(u)&=\frac1{1+u},\\
 A_2(u)&=\frac{1-2u^2}{2(1+u)^3},\\
 A_3(u)&=\frac{(u-2)(6u^3+6u^2-1)}{6(1+u)^5},\\
 A_4(u)&=-\frac{24u^6-72u^5-96u^4+40u^3+69u^2+8u-6}
 {24(1+u)^7},\\
 A_5(u)&=\frac{120u^8-720u^7-480u^6+1380u^5+1320u^4
 -151u^3-428u^2-58u+24}
 {120(1+u)^9}.
\end{align*}
The limits \(A_n(0)=1/n\) are transparent in this coordinate.

\section[Derivation of the joint (q,1/L) expansion]{Derivation of the joint \((q,1/L)\) expansion}

This appendix records the calculation behind
\eqref{eq:joint-resummation}.  Set
\begin{equation}\label{eq:appendix-v-u-ansatz}
  v=t+uv_1+u^2v_2+u^3v_3+\cdots,
  \qquad t=-\log(1-q),\quad s=1-q.
\end{equation}
Since \(\e^{-t}=s\), substitution into
\[
  q=\frac1{1-uv}-\e^{-v}
\]
and comparison of powers of \(u\) gives
\begin{align}
 v_1={}&-\frac{t}{s},
 \label{eq:v1-joint}\\
 v_2={}&\frac{t\bigl(2+(2q-1)t\bigr)}{2s^2},
 \label{eq:v2-joint}\\
 v_3={}&-\frac{t}{6s^3}
 \left[
 6+(18q-9)t+(2-6q+6q^2)t^2
 \right].
 \label{eq:v3-joint}
\end{align}
Therefore
\begin{align}
 V(q,u)={}&t-\frac{ut}{s}
 +\frac{u^2t\bigl(2+(2q-1)t\bigr)}{2s^2}
 \notag\\
 &-\frac{u^3t}{6s^3}
 \left[6+(18q-9)t+(2-6q+6q^2)t^2\right]
 +\Ord(u^4).
 \label{eq:joint-through-u3}
\end{align}
At \(q=0\), every term vanishes, as required.  Expanding
\eqref{eq:joint-through-u3} again in powers of \(q\) reproduces the
large-\(L\) expansions of all \(\Acoef_n(L)\) simultaneously.

\section{Stirling polynomials for the Lambert expansion}

The polynomials \(P_n\) defined in \eqref{eq:Pn-def} through order eight
are
\begin{align*}
 P_1(s)={}&s,\\
 P_2(s)={}&\frac{s(s-2)}2,\\
 P_3(s)={}&\frac{s(2s^2-9s+6)}6,\\
 P_4(s)={}&\frac{s(3s^3-22s^2+36s-12)}{12},\\
 P_5(s)={}&\frac{s(12s^4-125s^3+350s^2-300s+60)}{60},\\
 P_6(s)={}&\frac{s(20s^5-274s^4+1125s^3-1700s^2+900s-120)}{120},\\
 P_7(s)={}&\frac{s}{840}
 \bigl(120s^6-2058s^5+11368s^4-25725s^3\\
 &\hspace{28mm}{}+24500s^2-8820s+840\bigr),\\
 P_8(s)={}&\frac{s}{2520}
 \bigl(315s^7-6534s^6+45962s^5-142149s^4\\
 &\hspace{28mm}{}+205800s^3-135240s^2+35280s-2520\bigr).
\end{align*}
They provide an arbitrary-order route from the exact \(W\)-coordinate
coefficients \(\Ccoef_n(W(z))\) to polynomials in
\(\log z\) and \(\log\log z\).

\section{A compact proof of the leading power-log inverse}

For completeness, we verify \eqref{eq:power-log-leading-inverse}.  Let
\[
  Y=cX^a(\log X)^b,
  \qquad b\ne0,
\]
and put \(s=\log X\).  Then
\[
  \left(\frac{Y}{c}\right)^{1/b}=s\e^{(a/b)s}.
\]
Multiplication by \(a/b\) gives
\[
  \frac ab\left(\frac{Y}{c}\right)^{1/b}
  =\frac ab s\,\exp\!\left(\frac ab s\right).
\]
Applying a suitable branch of \(W\),
\[
  \frac ab s=W_k\!\left(
  \frac ab\left(\frac{Y}{c}\right)^{1/b}\right),
\]
so \(X=\e^s\) is exactly \eqref{eq:power-log-leading-inverse}.  This
calculation illustrates why Lambert \(W\) is a natural normalizer for
series reversal with a power-logarithmic leading monomial.

\section{Reference identities and sanity checks}

The following identities are useful for independent verification:
\begin{align}
 &L\e^L=z,
 &&\e^{-L}=\frac{L}{z},
 &&L'=\frac{L}{z(1+L)},
 \label{eq:check-L}\\
 &\rW_1(z)\e^{\rW_1(z)}+\rW_1(z)=z,
 &&g(z)=z-\rW_1(z),
 &&W(g(z))=\rW_1(z),
 \label{eq:check-rW}\\
 &g(z)+W(g(z))=z,
 &&g'(z)=\frac1{1+W'(g(z))},
 &&z-W(z)\le g(z)\le z.
 \label{eq:check-g}
\end{align}
For symbolic coefficients, three equivalent checks are available:
\begin{enumerate}
\item compare \(\Acoef_n\) from ordinary Lagrange inversion with
      \(\Ccoef_n/L^n\) from the differential operator;
\item substitute the truncated \(v\)-series into
      \(B_L(v)-q\) and verify cancellation;
\item substitute \(G_N\) into \(G_N+W(G_N)-z\) and verify the predicted
      residual order.
\end{enumerate}

\begin{thebibliography}{99}

\bibitem{Corless1996}
R.~M. Corless, G.~H. Gonnet, D.~E.~G. Hare, D.~J. Jeffrey, and D.~E. Knuth,
\newblock On the Lambert \(W\) function,
\newblock \emph{Advances in Computational Mathematics} \textbf{5} (1996),
329--359.

\bibitem{DLMF}
NIST Digital Library of Mathematical Functions,
\newblock \S4.13, Lambert \(W\)-function,
\newblock F.~W.~J. Olver et al., editors,
\newblock \url{https://dlmf.nist.gov/4.13}.

\bibitem{MezoBaricz2017}
I.~Mez\H{o} and \'A.~Baricz,
\newblock On the generalization of the Lambert \(W\) function,
\newblock \emph{Transactions of the American Mathematical Society}
\textbf{369} (2017), no.~11, 7917--7934,
\newblock DOI: \href{https://doi.org/10.1090/tran/6911}{10.1090/tran/6911}.

\bibitem{Comtet}
L.~Comtet,
\newblock \emph{Advanced Combinatorics: The Art of Finite and Infinite
Expansions},
\newblock D. Reidel, Dordrecht, 1974.

\bibitem{FlajoletSedgewick}
P.~Flajolet and R.~Sedgewick,
\newblock \emph{Analytic Combinatorics},
\newblock Cambridge University Press, Cambridge, 2009.

\bibitem{Henrici}
P.~Henrici,
\newblock \emph{Applied and Computational Complex Analysis}, Vol.~2,
\newblock Wiley, New York, 1977.

\bibitem{Olver}
F.~W.~J. Olver,
\newblock \emph{Asymptotics and Special Functions},
\newblock A K Peters, Wellesley, MA, 1997.

\bibitem{vanDerHoeven}
J.~van der Hoeven,
\newblock \emph{Transseries and Real Differential Algebra},
\newblock Lecture Notes in Mathematics 1888, Springer, Berlin, 2006.

\bibitem{ProveItTransseries}
V.~Reshetnikov,
\newblock Series and transseries draft packages,
\newblock ProveIt repository, polynomial-logarithmic transseries and
transseries tutorials,
\newblock \href{https://github.com/VladimirReshetnikov/ProveIt/tree/main/Analysis/FabiusFunction/docs/semi-formalized-research-frontiers/drafts/series-and-transseries}{Repository directory: \texttt{series-and-transseries}}.

\end{thebibliography}

\end{document}
